\introduction % Структурный элемент: ВВЕДЕНИЕ

В современном цифровом мире, где данные стали одним из ключевых активов, эффективность их обработки и хранения напрямую определяет успешность бизнес-процессов, качество пользовательского опыта и, в конечном счёте, финансовые результаты организаций. Основой большинства информационных систем остаются системы управления базами данных, от производительности которых зависит скорость отклика приложений, стабильность работы под нагрузкой и возможность масштабирования. Особую значимость эти вопросы приобретают в контексте постоянного роста объёмов данных, усложнения запросов и повышения требований к доступности сервисов.

Однако существующие инструменты для нагрузочного тестирования баз данных зачастую не предоставляют комплексного решения. Многие из них требуют глубоких технических знаний для настройки, проводят тестирование в изолированной среде без возможности простого сопоставления результатов нескольких прогонов, а их выводы ограничиваются текстовыми отчётами или сырыми данными, требующими трудоемкой ручной обработки и анализа. Отсутствие интегрированной платформы, которая бы автоматизировала процесс тестирования, обеспечивала наглядную визуализацию ключевых метрик в реальном времени и позволяла проводить удобное сравнительное исследование, создаёт существенный барьер для оперативной диагностики узких мест и обоснованного принятия решений по оптимизации. В связи с этой проблемой становится актуальной разработка программного комплекса, который заполнит указанный пробел.

Целью данной выпускной квалификационной работы является разработка и внедрение программного комплекса, позволяющего в автоматическом режиме проводить нагрузочное тестирование реляционных СУБД, собирать ключевые метрики производительности и предоставлять результаты в наглядном виде через единый веб-интерфейс для последующего анализа и сравнения.

Для достижения поставленной цели необходимо решить следующие задачи:
\begin{enumerate}[label=\arabic*)]
    \item провести анализ современных подходов к нагрузочному тестированию СУБД и сформулировать требования к разрабатываемой системе;
    \item разработать и реализовать ядро системы для автоматизированного тестирования производительности с поддержкой SQL-сценариев и транзакционной нагрузки;
    \item спроектировать и создать веб-интерфейс для управления процессом тестирования и визуального анализа результатов;
    \item выполнить комплексное тестирование системы на тестовых стендах и проанализировать полученные результаты.
\end{enumerate}

Объектом исследования выступают реляционные системы управления базами данных и их поведение под контролируемой нагрузкой. Предметом исследования являются методы, алгоритмы и инструментальные средства для автоматизированного проведения нагрузочного тестирования, сбора метрик производительности и визуализации результатов с целью сравнительного анализа.

В работе применяются методы системного анализа для проектирования архитектуры, методы объектно-ориентированного проектирования для реализации модулей системы, методы математической статистики для сравнительного анализа результатов тестирования, а также методы экспериментального исследования для оценки производительности СУБД.

Практическая значимость работы заключается в создании инструмента, который позволяет не только измерять, но и наглядно сравнивать производительность различных СУБД, что является основой для обоснованного выбора технологий, оценки эффективности оптимизаций и предотвращения регрессии после внесения изменений. Целевой аудиторией системы являются инженеры по производительности, DevOps-инженеры, администраторы баз данных и разработчики.

Структура работы включает введение, четыре главы, заключение и список использованных источников. В первой главе проводится аналитический обзор существующих решений и формулируются требования. Во второй главе описывается проектирование архитектуры и модулей системы. В третьей главе рассматривается программная реализация. В четвёртой главе представлены результаты экспериментального тестирования~--- двенадцать прогонов~--- и анализ производительности на двух тестовых наборах данных.